\section{Discussion}\label{discussion}

This model offers a radical reinterpretation of core physical concepts
by shifting the ontological basis of physics from abstract mathematical
fields to a concrete, tensioned medium --- a 3D brane embedded in 4D
space. By doing so, it challenges the standard partitioning of physics
into separate frameworks for general relativity, quantum mechanics, and
field theory, and instead suggests that all three may be different
aspects of wave propagation and internal tension dynamics within a
unified medium.

\paragraph{Relation to Randall--Sundrum brane models.}
The present use of a ``brane'' as a physical elastic medium embedded in four-dimensional Euclidean space should not be confused with Randall--Sundrum--type brane-world scenarios~\cite{RandallSundrum1999}. In Randall--Sundrum models, a thin 3+1-dimensional brane with tuned tension is embedded in an AdS$_5$ bulk, and four-dimensional gravity arises from a localized graviton zero mode of a fundamental five-dimensional Einstein theory. By contrast, the current framework postulates no higher-dimensional Einstein--Hilbert action: the only fundamental degree of freedom is the elastic embedding field $\mathbf{X}(x,t)$, the effective metric seen by internal waves is the induced metric on the brane, and gravity is interpreted as an emergent manifestation of tension-driven lateral contraction fields sourced by internal energy concentrations in the medium.

The strength of this approach lies in its coherence and minimalism:
rather than postulating multiple fundamental fields, point particles, or
probabilistic axioms, it reduces the complexity to a single substrate
with consistent rules. Properties like mass, spin, and even gauge
symmetry emerge not through external prescriptions but through the
solitonic geometry and oscillatory modes of the medium. Gravitation
becomes a local contraction effect tied to transverse deformation, while
magnetic fields arise from relativistic reinterpretations of electric
fields --- reinforcing the idea that secondary forces may reflect
higher-order dynamical effects of a more fundamental structure.

Several parallels with analog gravity systems~\cite{Barcelo2011AnalogueGravity,Volovik2003HeliumDroplet}, such as Bose-Einstein
condensates and fluid simulations, suggest that emergent relativity from
wave propagation is not a mere metaphor, but a physically plausible
mechanism. Furthermore, the topological analogies to Skyrmions and Hopf
solitons provide grounding for the idea that stable particles can arise
as localized, non-dissipative excitations with internal symmetry and
spin structure.

\paragraph{Robustness of parameter identifications.}
Among the various parameter identifications, the mapping between the brane wave speed and the speed of light $c$, and the use of the Compton frequency to define $\hbar$, are relatively robust consequences of the elastic brane dynamics. By contrast, the proposed mappings for the elementary charge $e$ and the gravitational constant $G$ involve more phenomenological assumptions about how internal observers infer effective electromagnetic and gravitational fields. These identifications should therefore be regarded as conjectural and subject to further refinement.

We also note that the de~Broglie--Einstein ``harmony of phases'' is
naturally realized in this framework. The Compton-frequency internal
mode of the toroidal electron, modeled as a lightlike amplitude wave
$\xi(t,z,x,y)$ in a curved world-tube, propagates at the microscopic
wave speed $c$ in the effective brane metric, while the soliton
center-of-mass follows a timelike trajectory. As shown in
Section~\ref{subsec:harmony-phases}, this geometric setup
automatically leads to the usual de~Broglie phase relation
$\phi(t,x) = \omega_C \gamma (t - v x/c^2)$ and the corresponding
de~Broglie wavelength when expressed in lab coordinates, without
postulating Lorentz symmetry at the microscopic level.

\paragraph{Dirac equation as emergent envelope dynamics.}

The reinterpretation of the Dirac equation adopted here should be
emphasized. We do not propose to modify the empirical content of the
Dirac theory; instead, we embed it into a more microscopic picture in
which the electron is realized as a toroidal brane soliton and the
relativistic Dirac spinor $\psi(x)$ provides an efficient description
of the soliton's centre-of-mass and spin degrees of freedom in the
long-wavelength limit. The microscopic dynamics remain those of the
nonlinear brane, with emergent Lorentz symmetry arising from the
tension and lateral excitation structure as discussed in
Section~\ref{sec:emergent-relativity}. The Dirac current and spin
bilinears then encode, at a coarse-grained level, the soliton
worldline and the orientation and handedness of the toroidal internal
mode, while the de~Broglie phase of $\psi(x)$ is identified with the
phase of the internal Compton clock in the tubular description of
Section~\ref{subsec:tubular-electron}. In this sense the Dirac
equation is not a competing microscopic postulate, but an emergent
envelope equation obeyed by the effective electron degrees of freedom
in the brane model.

From this vantage point, matter and radiation differ only in how the
underlying energy is organized: solitons correspond to ``canned''
energy, confined and stabilized by the brane's tension and nonlinear
self-interaction, whereas free photon modes correspond to ``released''
energy propagating as nearly linear waves. This echoes van der Mark and
't~Hooft's characterization of matter as ``canned energy'' and radiation
as ``free energy'' in their analysis of mass and light
\cite{vanderMark2000LightHeavy}, but here that slogan is given a concrete
mechanical realization in terms of explicit brane deformations.

\paragraph{Geometric phase readout (Berry) and its current scope.}
A central motivation of this project is that a purely mechanical substrate may encode nontrivial phase
transport in the \emph{direction} of a prepared narrowband excitation, even when the underlying fields
$\mathbf X(x,t)$ and $\partial_t\mathbf X(x,t)$ are real.
In the present revision we take a deliberately conservative step: we define only the minimal $U(1)$
Berry connection/phase readout on normalized complex envelopes (Section~\ref{sec:brane-to-berry-to-em})
and restrict attention to robust loop phases (holonomies) computed from overlaps.

This step already imposes a concrete requirement: the prepared packet must occupy an effectively isolated
branch on its spectral support (a practical gap condition). Establishing such an isolated band is therefore
a parameter-and-initialization problem first, and only afterwards an analysis problem.
More ambitious claims --- mapping the Berry connection to an effective electromagnetic four-potential,
reconstructing curvature/source distributions, or invoking non-Abelian ($U(N)$) transport near degeneracies ---
are postponed until the minimal holonomy measurements are numerically stable and reproducible.

\paragraph{Curvature toy models as secondary diagnostics.}
Toy calculations that translate localized 4D bulges into intrinsic curvature measures can serve as
sanity checks for the emergent-gravity intuition, but they are not central to the present validation
path. The current priority is the initialization targets and the Berry/holonomy pipeline; curvature
estimates will be revisited once the core diagnostics are stable and reproducible.

While not yet a full theory, this work constitutes a philosophical and
mathematical prototype --- a sandbox in which gravity, quantum behavior,
and particle identity emerge from the interplay of wave dynamics,
tension, and geometry. Its value may lie as much in the questions it
inspires as in the answers it attempts to provide. Further refinement
and critical engagement from both theorists and computational physicists
may help clarify whether this direction offers a fruitful path toward a
more intuitive and unified understanding of the physical world.